\begin{lemma}
\label{lemma-2-3-ring-map}
Let $\varphi : R \to S$ be a ring map. Assume
\begin{enumerate}
\item[(a)] $S$ is generated as an $R$-algebra by elements $x$ such
that $x^2, x^3 \in \varphi(R)$, and
\item[(b)] $\Ker(\varphi)$ is locally nilpotent,
\end{enumerate}
Then $\varphi$ induces isomorphisms on residue fields and
a homeomorphism of spectra. For any ring map $R \to R'$
the ring map $R' \to R' \otimes_R S$ also satisfies (a) and (b).
\end{lemma}

\begin{proof}
For every original generator $x$, choose $a_x,b_x\in R$
with $\varphi(a_x)=x^2$ and $\varphi(b_x)=x^3$.
The monic equation $T^2-a_x$ proves integrality of
$x$ over $R$. Every element of the generated algebra
involves finitely many such generators, so the
integral-closure subring property proves that
$R\to S$ is integral. Put $I=\ker\varphi$.
The integral injection $R/I\to S$ is surjective
on spectra by Lemma \ref{lemma-integral-overring-surjective}.
Since $I$ is locally nilpotent, $V(I)=\Spec(R)$.
Thus the original map on spectra is surjective,
and is closed by Lemmas \ref{lemma-integral-going-up}
and \ref{lemma-going-up-closed}.

Let $\mathfrak q$ lie over $\mathfrak p$, and use
the original injection
$\kappa(\mathfrak p)\to\kappa(\mathfrak q)$.
If $\alpha$ is the image of $x$ in the receiving
field, then $\alpha^2=\overline a_x$ and
$\alpha^3=\overline b_x$, where the bars are the
images of these original coefficients in
$\kappa(\mathfrak p)$. If $\overline a_x=0$,
then $\alpha=0$. Otherwise
$\alpha=\overline b_x/\overline a_x$.
Consequently every generator belongs to the
image of the original base residue field.
Since $\kappa(\mathfrak q)=\operatorname{Frac}(S/\mathfrak q)$
is generated as a field by those images and
$\kappa(\mathfrak p)$, this residue-field map
is an isomorphism.

If $\mathfrak q'$ also lies over $\mathfrak p$,
compose the maps $S\to\kappa(\mathfrak q)$ and
$S\to\kappa(\mathfrak q')$ with the respective
inverses of the base residue-field isomorphisms.
Both resulting maps to $\kappa(\mathfrak p)$
agree on $R$ and send each original $x$ to
the same displayed value, zero or
$\overline b_x/\overline a_x$.
They agree on every polynomial in the
generators, hence on $S$. Their kernels
are respectively $\mathfrak q$ and
$\mathfrak q'$, so these primes are equal.
The original spectrum map is a continuous
closed bijection, hence a homeomorphism.

Retain any original $f:R\to R'$ and put
$S'=R'\otimes_R S$. It is generated over
$R'$ by $1\otimes x$, with
$$
(1\otimes x)^2=f(a_x)\otimes1,
\qquad (1\otimes x)^3=f(b_x)\otimes1.
$$
To verify the new kernel without assuming
$f$ flat, let $\mathfrak p'\subset R'$
contract to $\mathfrak p\subset R$.
Choose an original $\mathfrak q$ above
$\mathfrak p$, as is possible by surjectivity.
The tensor ring
$$
C=\kappa(\mathfrak p')\otimes_{\kappa(\mathfrak p)}
\kappa(\mathfrak q)
$$
is nonzero: a vector-space basis of either
nonzero field remains a basis after extending
scalars. The original maps give
$S'\to C$, $r'\otimes s\mapsto\overline r'\otimes\overline s$.
Take a prime of $C$. Its quotient is a
nonzero domain, and the induced unital
map from the field $\kappa(\mathfrak p')$
is injective. The inverse image of that
prime in $S'$ therefore contracts to
exactly $\mathfrak p'$ in $R'$.
Thus $\Spec(S')\to\Spec(R')$ is surjective,
and its kernel lies in every prime of
$R'$, hence consists of nilpotents.
This proves both base-change hypotheses;
the already proved argument gives the
same homeomorphism and residue-field
isomorphisms after every such base change.

The exponents $2,3$ can more generally
be replaced, separately for each original
generator $x$, by a finite nonempty list
of positive integers $d_1,\ldots,d_r$
with greatest common divisor $1$, provided
all $x^{d_i}$ belong to $\varphi(R)$.
One positive exponent supplies the monic
integral equation. Choose integers $u_i$
with $\sum_i u_i d_i=1$. For a nonzero
residue $\alpha$ of $x$, every $\alpha^{d_i}$
is a nonzero element of the base residue
field, and the full identity
$$
\alpha=\prod_{i=1}^r(\alpha^{d_i})^{u_i}
$$
puts it in that field, with negative
powers interpreted through its actual
field inverses. If any power is zero,
then $\alpha=0$. This formula determines
the same residue at every prime over
the given base prime. The preceding
integrality, kernel, and generator-map
arguments therefore prove the same
conclusions and their preservation by
arbitrary base change for these lists.

The greatest-common-divisor condition
cannot be weakened to an arbitrary
positive common divisor. For a list
with greatest common divisor $d>1$,
take the original field inclusion
$\mathbf Q(u^d)\to\mathbf Q(u)$, with
$u$ transcendental. The powers in the
list all belong to the smaller field.
The elements $1,u,\ldots,u^{d-1}$ are
independent over $\mathbf Q(u^d)$:
clearing the original denominators
leaves polynomial monomials with
distinct exponent residues modulo $d$.
They span the algebra generated by $u$
using $u^d\in\mathbf Q(u^d)$; that
finite-dimensional domain is a field,
since multiplication by a nonzero
element is injective and hence
surjective on its finite-dimensional
vector space. Thus it is the original
$\mathbf Q(u)$ and the residue-field
extension has degree exactly $d$,
not degree one.

Even under the original $2,3$ hypotheses,
an injective map can acquire a nonzero
nilpotent kernel after base change.
For any field $k$ let
$$
R=k[a,b]/(b^2-a^3)\longrightarrow S=k[t],
\qquad a\longmapsto t^2,\quad b\longmapsto t^3.
$$
This is injective: monic division in
$b$ writes an element uniquely as
$f(a)+b g(a)$, and its image
$f(t^2)+t^3g(t^2)$ has disjoint even
and odd exponent supports. Also $S$
is generated by the original $t$ with
$t^2,t^3$ in the image. Base change
by $R\to R/(a)=k[b]/(b^2)$ gives
$$
k[b]/(b^2)\longrightarrow k[t]/(t^2),
\qquad b\longmapsto t^3=0.
$$
Its kernel is exactly the nonzero
ideal $(b)$, whose square is zero.
The quotient-tensor comparison
$(R/(a))\otimes_R S\simeq S/(t^2)$
sends $[r]\otimes s$ to $[\varphi(r)s]$
and has inverse $[s]\mapsto1\otimes s$.
Thus this is the actual base-changed
map, with its nilpotent kernel fully
identified.
\end{proof}